\Article{39}
Der Wert von $v$, der im Allgemeinen eine Funktion von $r$, $u$ und $\lambda$ ist, geht auf der Erdoberfläche in eine Funktion allein von $u$ und $\lambda$ über, und
\[
-\frac{\mathrm dv}{R\mathrm du},\qquad -\frac{\mathrm dv}{R\sin u\,\mathrm d\lambda}
\]
sind die horizontalen Komponenten der dort durch die galvanischen Ströme hervorgerufenen magnetischen Kraft, die nach Norden beziehungsweise nach Westen gerichtet sind. Offenbar gelten also die bemerkenswerten, oben in Artikel 15 und 16 angeführten Sätze auch hier. Bei der dritten Komponente, der vertikalen magnetischen Kraft, verhält es sich jedoch etwas anders, wenn die Ursachen oberhalb der Erde liegen, als wenn sie sich in ihrem Inneren befinden. Um die aus den oberhalb liegenden Ursachen entstehende vertikale Kraft zu bestimmen, muss man $v$ zunächst als Funktion aller drei Variablen $r$, $u$ und $\lambda$ betrachten, nach $r$ differenzieren und dann $r=R$ einsetzen. Für den inneren Raum $S'$, zu dem die Erdoberfläche gehört, lässt sich $v$ aber nur in eine Reihe nach steigenden Potenzen von $r$ entwickeln. Setzen wir
\[
\frac vR=p^0+\frac rR\cdot p'+\frac{rr}{RR}\cdot p''+\frac{r^3}{R^3}\cdot p'''+\text{u. s. w.}
\]
so ist $p^0$ eine Konstante, nämlich der Wert von $\frac vR$ im Mittelpunkt der
\SourceLeaf{176}
Erde. Dagegen sind $p'$, $p''$, $p'''$ usw. Funktionen von $u$ und $\lambda$, die dieselben partiellen Differentialgleichungen erfüllen wie oben $P'$, $P''$, $P'''$ usw. Daraus folgt ähnlich wie oben in Artikel 20, dass es genügt, den Wert von $v$ in jedem Punkt der Erdoberfläche zu kennen, um daraus den allgemeinen, im ganzen Raum $S'$ gültigen Ausdruck abzuleiten. Den Wert bis auf einen konstanten Anteil oder, was dasselbe ist, die Koeffizienten $p'$, $p''$, $p'''$ usw. kann man bereits aus den horizontalen Kräften auf der Erdoberfläche bestimmen. Der Wert der vertikalen Kraft auf derselben Oberfläche ist aber nicht
\[
=2p'+3p''+4p'''+\text{u. s. w.}
\]
(wie er es wäre, wenn die Kräfte aus dem Erdinneren wirkten). Er ist vielmehr
\[
=-p'-2p''-3p'''-\text{u. s. w.}
\]

Da nun unsere numerischen Elemente (Artikel 26), die unter der Voraussetzung der ersten Formel bestimmt wurden, die Gesamtheit der Erscheinungen bereits sehr zufriedenstellend darstellen, während diese mit der zweiten Formel völlig unvereinbar wären, ist die Unhaltbarkeit der Hypothese, die die Ursachen des Erdmagnetismus in den Raum außerhalb der Erde verlegt, als erwiesen anzusehen.

\Article{40}
Damit darf jedoch die Möglichkeit, dass ein \textit{Teil} der erdmagnetischen Kraft, wenn auch nur ein vergleichsweise sehr geringer, von oben her erzeugt wird, noch nicht als eindeutig widerlegt gelten. Eine viel vollständigere und viel genauere Kenntnis der Erscheinungen wird künftig über diesen wichtigen Punkt der Theorie Aufschluss geben. Geht man von gemischten Ursachen aus und versteht die Zeichen $V$, $P^0$, $P'$, $P''$ usw., $v$, $p^0$, $p'$, $p''$ in derselben Bedeutung wie oben, sodass sich die ersteren auf die aus dem Inneren und die letzteren auf die aus dem Außenraum wirkenden Ursachen beziehen, und setzt man ferner
\[
V+v=W,\quad P^0+p^0=\Pi^0,\quad P'+p'=\Pi',\quad P''+p''=\Pi''\quad\text{u. s. w.}
\]
\SourceLeaf{177}
so gilt auf der Erdoberfläche
\[
\frac WR=\Pi^0+\Pi'+\Pi''\quad\text{u. s. w.}
\]
wobei $\Pi^{\mathrm n}$ dieselbe partielle Differentialgleichung erfüllt wie $P^{\mathrm n}$ (Artikel 18). Die beiden Komponenten der dort wirkenden horizontalen magnetischen Kraft werden durch
\[
-\frac{\mathrm dW}{R\mathrm du},\qquad-\frac{\mathrm dW}{R\sin u\,\mathrm d\lambda}
\]
ausgedrückt. Die in Artikel 15 und 16 angeführten Sätze bleiben also auch hier gültig. Aus der bloßen Kenntnis der horizontalen Kräfte kann man die Größen $\Pi'$, $\Pi''$, $\Pi'''$ usw. bestimmen, daraus allein aber gar nichts über das Vorliegen gemischter Ursachen schließen. Betrachtet man jedoch die vertikale Kraft für sich und bringt sie in die Form
\[
Q^0+Q'+Q''+Q'''+\text{u. s. w.}
\]
sodass $Q^{\mathrm n}$ die zuvor erwähnte partielle Differentialgleichung erfüllt. so gilt
\[
\begin{gathered}
Q^0=P^0,\quad Q'=2P'-p',\quad Q''=3P''-2p'',\quad Q'''=4P'''-3p'''
\end{gathered}
\]
usw., und folglich
\[
\begin{aligned}
3P'&=\Pi'+Q',&3p'&=2\Pi'-Q'\\
5P''&=2\Pi''+Q'',&5p''&=3\Pi''-Q''\\
7P'''&=3\Pi'''+Q''',&7p'''&=4\Pi'''-Q'''\quad\text{u. s. w.}
\end{aligned}
\]

Verbindet man also die horizontalen Kräfte mit der vertikalen Kraft, so kann man $W$ in seine Bestandteile $V$ und $v$ zerlegen und dadurch erkennen, ob dem letzteren ein merklicher Wert zugeschrieben werden muss. Nur den konstanten Anteil von $v$, nämlich $p^0$, lassen die Beobachtungen völlig unbestimmt. Der Grund dafür ergibt sich aus Artikel 38.

Auch aus diesem interessanten Blickwinkel erscheint es daher wichtig, die horizontale magnetische Kraft für sich zu betrachten. Darin sehen wir einen weiteren Grund, die oben in Artikel 21 empfohlenen Gesichtspunkte zu berücksichtigen.

\SourceLeaf{178}
\Article{41}
Für die im vorigen Artikel angedeutete Untersuchung wird es wahrscheinlich noch lange an ausreichenden Daten fehlen. Es ist aber bemerkenswert, dass sich die Schwankungen der magnetischen Kraft. wie sie sich gleichzeitig an den verschiedenen Punkten der Erdoberfläche zeigen, ganz ähnlich behandeln lassen. Dafür lassen sich vielleicht schon sehr viel früher die unbedingt nötigen Daten zusammentragen. Dies gilt sowohl für die regelmäßigen, mit der Tages- und Jahreszeit wechselnden Änderungen als auch für die unregelmäßigen. Einige allgemeine Hinweise zu diesen künftigen Untersuchungen dürfen hier wohl noch Platz finden.

Nachdem man die beobachteten gleichzeitigen Änderungen für jeden Ort als Änderungen der Komponenten der magnetischen Kraft, $\Delta X$, $\Delta Y$, $\Delta Z$, ausgedrückt hat, muss man zuerst untersuchen, ob sich die Änderungen der beiden horizontalen Komponenten gemäß unserer Theorie verhalten. Ihr zufolge müssen $-\Delta X$ und $-\sin u.\Delta Y$ die Werte der partiellen Ableitungen einer Funktion von $u$ und $\lambda$ nach diesen Variablen sein. Ist das der Fall, so wird man schließen, dass die Ursachen entweder wirkliche galvanische Ströme sind oder doch wenigstens so wirken wie diese oder wie getrennte magnetische Flüssigkeiten. Im entgegengesetzten Fall wäre erwiesen, dass die Ursachen keine galvanischen Ströme sein können. Man sieht, dass schon die Kenntnis solcher Änderungen der horizontalen Kraft allein höchst wichtige Aufschlüsse geben kann, sofern die Beobachtungen hinreichend genau, zahlreich und räumlich verteilt sind. Kennt man darüber hinaus die gleichzeitigen Änderungen der vertikalen Kraft, so wird die Methode des vorigen Artikels \textit{unter der Voraussetzung des ersten Falls} Auskunft darüber geben, ob die Ursachen oberhalb oder unterhalb der Erdoberfläche liegen. Da sich diese Ursachen wahrscheinlich in einer im Vergleich zum ganzen Erdkörper recht dünnen Schicht befinden, wird sich dann sogar ihre räumliche Verteilung wenigstens näherungsweise bestimmen lassen.

Was dagegen den oben als möglich erwähnten zweiten Fall betrifft, so glaube ich ihn für die regelmäßigen, von der Tages- und Jahreszeit abhängigen Änderungen der erdmagnetischen Kraft für wenig wahrscheinlich halten zu dürfen. Bei den unregelmäßigen, in kurzen Zeitabständen wechselnden Änderungen würde ich zurzeit aber kaum wagen, in \SourceLeaf{179}dieser Hinsicht eine Vermutung zu äußern. Sollten sie ihren Ursprung in großen Bewegungen der Elektrizität oberhalb der Atmosphäre haben, so ließen sich diese schwerlich zu den galvanischen Strömen zählen. Denn obwohl alles dafür spricht, galvanischen Strom als Elektrizität in Bewegung anzusehen, ist eine Bewegung der Elektrizität nur dann ein galvanischer Strom, wenn sie einen in sich zurückkehrenden Kreislauf bildet. Da sich nur unter dieser Bedingung, wie mehrfach erwähnt, getrennte magnetische Flüssigkeiten an die Stelle des galvanischen Stroms setzen lassen, würden unsere Beziehungen zwischen den Komponenten unter der genannten Hypothese nicht mehr gelten. Das heißt, der zweite Fall würde tatsächlich eintreten. Diese wichtige Tatsache mit Gewissheit festzustellen, wäre schon an sich von großem Interesse. Doch selbst dann könnte es mit hinreichend weit verteilten und zuverlässigen Beobachtungen in unserer Reichweite liegen, herauszufinden, wo solche Bewegungen auftreten und wie sie sich verhalten.
